\section{行业前瞻：模型进步与 Agent 创新的边界}

\subsection{“模型会变强”与“Agent 仍有价值”并不矛盾}

季逸超反对“只要模型够强，Agent 层就没有价值”的观点。原因是现实任务依赖外部环境状态，而环境无法完整内化进参数。即使模型更强，任务编排、权限控制和执行验证仍然需要 Agent 层承担。

\begin{importantbox}{一个更稳妥的判断}
模型进步会降低 Agent 的部分实现成本，但不会消除 Agent 的系统职责。未来竞争更像“模型能力乘以系统能力”，而不是二选一。
\end{importantbox}

\subsection{多模态优先级：先输入理解，再输出花哨}

访谈中提到，很多场景真正痛点是“看懂复杂输入”，不是“生成漂亮输出”。例如截图、图表、半结构化文档、工具回传图片等，模型若在输入理解环节失真，后续链路很难纠正。

\begin{warningbox}{多模态常见误区}
把注意力集中在生成端（图像/视频输出）容易产生短期演示效果，但对 Agent 任务完成率贡献有限。对工作流产品而言，输入解析和状态跟踪往往更关键。
\end{warningbox}

\subsection{给创业者的建议：先定义问题，再匹配技术}

访谈里一条可执行建议是：先把问题定义清楚，再决定是做模型、做系统，还是做产品壳层。问题定义准确，技术路径才有收敛可能。否则团队会在“看起来都能做”的机会海里反复横跳。

\begin{knowledgebox}{问题定义模板（可直接套用）}
\begin{itemize}
    \item 用户目标：最终交付物是什么
    \item 约束条件：时间、权限、可审计性、安全边界
    \item 成功标准：一次完成率、重试成本、人工介入比例
\end{itemize}
\end{knowledgebox}

\subsection{本章小结}

对 Agent 创业来说，乐观要建立在结构化判断上：承认模型会持续进步，同时把系统层能力做到不可替代。


